\chapter{Planejamento do Estudo de Caso}
\label{cap:estudo_de_caso}

Este capítulo estabelece o planejamento do estudo de caso a ser conduzido na segunda etapa da pesquisa, formalizando os fundamentos metodológicos, o contexto do objeto investigado e os critérios que nortearão a coleta e a avaliação das evidências. Inicialmente, apresentam-se a definição do método empírico adotado e o objetivo específico do estudo. Na sequência, caracterizam-se o caso sob investigação, os trabalhos relacionados extraídos da revisão da literatura que embasam o protocolo e a questão de pesquisa do estudo de caso acompanhada das métricas normativas associadas.

\section{Definição}
\label{sec:ec_definicao}

Para \citeonline{yin2015}, o estudo de caso constitui uma investigação empírica voltada ao exame de um fenômeno contemporâneo em seu contexto real, particularmente adequada quando os limites entre o fenômeno observado e o ambiente que o envolve não podem ser separados com precisão. Na Engenharia de Software, \citeonline{runeson2009} reforçam a pertinência desse método porque os artefatos técnicos, os processos operacionais e as exigências organizacionais interagem de maneira complexa, impedindo que mecanismos de auditoria e segurança sejam avaliados isoladamente da infraestrutura em execução.

Ainda segundo \citeonline{yin2015} e \citeonline{runeson2009}, o estudo de caso caracteriza-se por um delineamento flexível e iterativo. Embora demande um protocolo rigoroso previamente estabelecido antes do trabalho de campo, sua estrutura permite refinar procedimentos de coleta e análise à medida que o pesquisador obtém compreensão mais profunda das particularidades técnicas do ambiente investigado, sem comprometer o rastreamento das decisões tomadas ao longo do processo.

Com o propósito de conferir robustez às conclusões e mitigar vieses decorrentes de uma única perspectiva, o estudo de caso apoia-se na triangulação de múltiplas fontes de evidência \cite{yin2015,runeson2009}. Quando registros provenientes de camadas distintas da infraestrutura convergem para o mesmo diagnóstico normativo, eleva-se a confiabilidade dos achados. Neste planejamento, essa convergência fundamenta o cruzamento entre arquivos declarativos de configuração, registros nativos do orquestrador, trilhas de auditoria do sistema hospedeiro e eventos centralizados na plataforma de SIEM.

Na taxonomia proposta por \citeonline{yin2015}, um estudo de caso único pode ser delineado de forma holística, quando examina o ambiente apenas em sua totalidade global, ou incorporada, quando investiga múltiplas subunidades de análise no interior de um mesmo caso. Esta pesquisa adota um estudo de caso único incorporado: o caso holístico corresponde à infraestrutura orquestrada por Kubernetes sob análise, ao passo que as subunidades de análise correspondem aos artefatos técnicos, manifestos de configuração, registros em tempo de execução e eventos correlacionados no SIEM avaliados sob os critérios do Controle 8 do PPSI 2.0 e da NC nº 21/GSI-PR. Esse arranjo permite examinar cada dimensão técnica em profundidade sem perder a visão integrada da conformidade global da infraestrutura.

Adicionalmente, \citeonline{runeson2009} enfatizam que as ameaças à validade devem ser consideradas desde a concepção do estudo, sob quatro dimensões clássicas igualmente adotadas por \citeonline{yin2015}. A validade de constructo verifica em que medida os indicadores operacionais representam fielmente o conceito de conformidade investigado pela questão de pesquisa. A validade interna examina riscos de fatores não controlados interferirem na atribuição dos efeitos observados entre a linha de base e o cenário pós-intervenção. A validade externa trata do alcance da generalização analítica dos achados para outras infraestruturas governamentais com características arquiteturais semelhantes. Por fim, a confiabilidade refere-se à reprodutibilidade dos procedimentos de coleta e análise, garantindo que outro pesquisador alcance conclusões equivalentes ao aplicar o mesmo protocolo.

\section{Objetivo}
\label{sec:ec_objetivo}

A verificação da segurança em infraestruturas computacionais exige critérios mensuráveis, objetivos e auditáveis que permitam diagnosticar lacunas reais e acompanhar a evolução de controles técnicos ao longo do tempo. Quando a conformidade regulatória é aferida por critérios subjetivos ou exclusivamente documentais, sem lastro em evidências extraídas do ambiente em operação, ampliam-se a vulnerabilidade operacional e o risco de descumprimento frente às auditorias dos órgãos de controle externo, problema agravado em arquiteturas dinâmicas baseadas em contêineres.

Assim, o objetivo deste estudo de caso consiste em investigar como uma infraestrutura orquestrada por Kubernetes combinada com uma plataforma de SIEM oferece mecanismos para verificar a conformidade perante as doze medidas do Controle 8 do PPSI 2.0 e os requisitos forenses da NC nº 21/IN01/DSIC/GSI-PR. Para tanto, o planejamento prevê o confronto sistemático entre o diagnóstico inicial de linha de base e o cenário posterior à intervenção arquitetural a ser conduzida na segunda etapa da pesquisa.

\section{Caso}
\label{sec:ec_caso}

O caso investigado neste trabalho corresponde à infraestrutura de computação em nuvem mantida pelo Laboratório Avançado de Produção, Pesquisa e Inovação em Software, conhecido como LabLivre/LAPPIS, unidade acadêmica de pesquisa e desenvolvimento vinculada à Universidade de Brasília (UnB). Por ser uma fundação pública federal integrante do Poder Executivo, a UnB compõe diretamente a Administração Pública Federal. Em decorrência desse vínculo institucional, os ambientes computacionais mantidos pelo laboratório, que abrangem tanto infraestruturas próprias de homologação, quanto plataformas digitais governamentais de relevância pública hospedadas em cooperação institucional, em tempo de desenvolvimento, como o GovHub e o Brasil Participativo \cite{brasilparticipativo2025}, submetem-se às diretrizes normativas expedidas pela Secretaria de Governo Digital e pelo Gabinete de Segurança Institucional da Presidência da República, em especial ao PPSI 2.0 \cite{portaria9511_2025,ppsi2026}, à Instrução Normativa SGD/MGI nº 4/2026 \cite{in4_2026} e à NC nº 21/GSI-PR \cite{nc21_2014}.

Sob o ponto de vista tecnológico, o ambiente sob análise adota o paradigma nativo em nuvem baseado na orquestração de contêineres por meio do Kubernetes \cite{cncf2024,kubernetes2025}. Mantendo independência em relação à topologia física específica ou a uma distribuição particular do orquestrador, a infraestrutura estrutura-se logicamente em duas camadas complementares: o plano de controle, responsável pelo gerenciamento declarativo do estado, pelo escalonamento de cargas e pelo processamento de requisições à API do orquestrador; e os nós de trabalho, encarregados da execução distribuída de aplicações \textit{web}, interfaces de programação de aplicações, sistemas gerenciadores de bancos de dados e serviços auxiliares de sustentação em contêineres de ciclo de vida efêmero.

A seleção desse caso fundamenta-se na combinação entre relevância institucional e viabilidade empírica de observação. As plataformas hospedadas no laboratório sustentam serviços públicos digitais e ciclos contínuos de desenvolvimento e evolução de software, reproduzindo fielmente os desafios de conformidade enfrentados por órgãos integrantes do Sistema de Administração dos Recursos de Tecnologia da Informação (SISP). Além disso, a acessibilidade operacional ao ambiente permite executar de maneira controlada e auditável as duas etapas previstas para a segunda fase da monografia, assegurando a coleta sistemática de artefatos antes e após a intervenção sem expor dados sensíveis ou comprometer a disponibilidade de serviços externos.

Este estudo de caso delimita-se às camadas de auditoria, registro de eventos e monitoramento de segurança em tempo de execução da infraestrutura orquestrada. Embora o ecossistema adote práticas modernas de entrega e operação de cargas em contêineres, infraestruturas dessa natureza apresentam como lacuna recorrente a ausência de centralização forense contínua sobre os eventos emitidos pelo plano de controle, pelos nós hospedeiros e pelos contêineres efêmeros, mantendo registros fragmentados localmente e desprovidos de correlação de segurança automatizada. O planejamento aqui definido endereça exatamente essa lacuna, estruturando a avaliação do estado inicial de conformidade e dos mecanismos técnicos necessários para suprir os requisitos normativos federais.

\todo[inline, color=pink]{Esta subseção precisa evoluir com mais detalhes o caso. Uma dos requisitos essenciais em um estudo de caso são os detalhes sobre o caso. Esses detalhes dizem respeito a quantidade de pessoas, perfil, principais tecnologias. Os projetos merecem um parágrado de explica
cão. Mas sobretudo e em especial, sobre a infraestrutura, incluindo topologia lógica, fluxos de dados, serviços críticos e quaisquer particularidades que possam impactar a coleta e análise de evidências.}

\section{Trabalhos Relacionados}
\label{sec:ec_trabalhos_relacionados}

Para sustentar o delineamento empírico, os mecanismos de coleta e a modelagem de conformidade deste estudo de caso, selecionaram-se cinco estudos primários identificados na Revisão Estruturada da Literatura apresentada no Capítulo~\ref{cap:revisao_sistematica}, os quais oferecem contribuições diretas e complementares ao protocolo proposto.

Para fundamentar o delineamento empírico, os mecanismos de coleta de dados e a modelagem de conformidade adotados neste estudo de caso, foram selecionados cinco estudos primários identificados na Revisão Estruturada da Literatura apresentada no Capítulo~\ref{cap:revisao_sistematica}. Esses estudos fornecem contribuições complementares que subsidiam a definição do protocolo proposto.


O trabalho de \citeonline{acosta2025} fundamenta a viabilidade de implantar o SIEM Wazuh em uma instituição pública real com restrições estritas de governança e soberania tecnológica, demonstrando o valor de uma etapa prévia de planejamento e diagnóstico de políticas institucionais antes da operacionalização do monitoramento centralizado. Embora o estudo valide o ganho operacional na detecção e no tratamento de incidentes em ambiente governamental, ele não mensura a evolução de conformidade por meio do catálogo de controle \textit{CIS} nem adota uma escala gradual padronizada, lacuna diretamente tratada neste estudo de caso.

A investigação de \citeonline{bungei2026} demonstra a aplicabilidade do SIEM Wazuh em ambientes com restrições orçamentárias e equipes enxutas, explorando a coleta de eventos em hospedeiros e o módulo de avaliação de configuração de segurança da plataforma para quantificar automaticamente a conformidade frente a catálogos \textit{CIS Benchmarks}. Essa abordagem embasa o uso de verificações automatizadas no SIEM como evidência objetiva de endurecimento; todavia, o estudo aplica a aferição apenas após a implantação em ambiente de testes, sem estabelecer uma linha de base normativa prévia nem abranger a auditoria do plano de controle Kubernetes.

Por sua vez, \citeonline{velasquez2025} abordam a estruturação de plataformas SIEM orientadas explicitamente à conformidade regulatória, propondo o mapeamento sistemático entre registros coletados, regras de análise e controles de referenciais internacionais como a ISO/IEC 27002 e a família de auditoria do NIST SP 800-53. Esse mapeamento reforça a premissa arquitetural desta pesquisa de vincular cada evento de auditoria e regra de relação, às exigências normativas rastreáveis, estendendo o conceito para o escopo das doze medidas do Controle 8 do PPSI 2.0 e da NC nº 21/GSI-PR.

No eixo de mensuração da maturidade, \citeonline{coutinho2023} destacam-se por aplicar uma escala ordinal formal de capacidade de processo baseada na norma ISO/IEC 15504 para avaliar controles de cibersegurança em uma infraestrutura operacional vinculada a instituição universitária. A proposta valida o rigor metodológico de traduzir evidências técnicas em níveis graduais de atendimento em vez de verificações puramente binárias, abordagem que nesta monografia é alinhada aos níveis oficiais de implementação e adoção estabelecidos pelo governo federal.

Por fim, \citeonline{torkura2021} estruturam um modelo de auditoria contínua de segurança em nuvem que exige a definição prévia de um estado inicial esperado dos recursos frente às políticas de conformidade antes do início do monitoramento contínuo de desvios. Essa lógica metodológica fundamenta a separação formal entre o diagnóstico inicial de linha de base e a avaliação pós-intervenção adotada no planejamento desta pesquisa, adaptando-a da auditoria via APIs de provedores públicos para a telemetria combinada entre Kubernetes e SIEM em infraestrutura governamental.

\section{Questão de Pesquisa}
\label{sec:ec_questao_pesquisa}

Seguindo o paradigma \textit{Goal Question Metric} (GQM) proposto por \citeonline{basili1994} e preservando o alinhamento vertical com a Questão Principal estabelecida na Seção~\ref{sec:intro_questao_pesquisa}, voltada a avaliar o grau de conformidade da infraestrutura de monitoramento de segurança em tempo de execução na Administração Pública Federal, formula-se a Questão de Pesquisa específica que norteia este estudo de caso:

\begin{quote}
\textbf{Como uma infraestrutura Kubernetes combinada com uma plataforma de SIEM pode oferecer mecanismos para verificar a conformidade da infraestrutura, segundo o Controle 8 do PPSI 2.0 e a NC 21/GSI-PR?}
\end{quote}

Para responder a essa questão de forma objetiva, associam-se exclusivamente os dois indicadores oficiais de maturidade definidos pela Instrução Normativa SGD/MGI nº 4/2026 \cite{in4_2026}, aplicados em dois momentos distintos da investigação, correspondentes ao diagnóstico inicial de linha de base e ao cenário posterior à implantação da arquitetura proposta:
\begin{itemize}
    \item \textbf{Nível de Implementação da Medida (NIM):} indicador ordinal aferido na escala gradual oficial de 0 a 5, com pontuações correspondentes de 0,00 a 1,00, para cada uma das doze medidas do Controle 8 do PPSI 2.0 \cite{ppsi2026,in4_2026}. A atribuição de cada nível ancora-se nos requisitos técnicos e forenses da NC nº 21/GSI-PR \cite{nc21_2014}, os quais abrangem sincronização temporal, metadados mínimos de eventos, centralização remota, integridade e retenção, atuando como critérios verificáveis sobre os artefatos coletados no diagnóstico inicial e no cenário pós-intervenção; e
    \item \textbf{Nível de Adoção do Controle (NAC):} indicador consolidado obtido pela média aritmética das pontuações NIM atribuídas às medidas aplicáveis do Controle 8, enquadrando o ambiente nas cinco faixas oficiais de adoção da IN SGD/MGI nº 4/2026 \cite{in4_2026} e permitindo quantificar de modo determinístico a evolução da conformidade entre a linha de base e o estado posterior à intervenção.
\end{itemize}

A adoção exclusiva do NIM e do NAC justifica-se por responder diretamente à lacuna identificada no Capítulo~\ref{cap:revisao_sistematica}, segundo a qual a literatura especializada carece de escalas graduais de maturidade alinhadas ao arcabouço regulatório nacional. Ao empregar a mesma métrica oficial exigida pela Secretaria de Governo Digital para o ciclo de governança dos órgãos do SISP, o protocolo assegura validade de constructo, comparabilidade auditável entre o estado inicial e o final e aplicabilidade imediata dos resultados na Administração Pública Federal.

\section{Fonte de Dados}
\label{sec:ec_fonte_dados}

Em consonância com a triangulação de evidências preconizada por \citeonline{yin2015} e \citeonline{runeson2009}, o protocolo organiza a coleta sobre quatro fontes complementares, integrando artefatos técnicos da infraestrutura em operação ao arcabouço normativo aplicável.

A primeira fonte compreende os arquivos declarativos e as configurações do ambiente orquestrado, abrangendo manifestos de política de auditoria do plano de controle, parâmetros de sincronização temporal nos nós hospedeiros e configurações de resolução de nomes e roteamento de entrada do \textit{cluster} Kubernetes \cite{kubernetes2025}. Esses artefatos comprovam se os mecanismos de registro estão habilitados e parametrizados conforme o padrão exigido.

A segunda fonte reúne os registros brutos de auditoria e sistema emitidos em tempo de execução, incluindo as trilhas do servidor de API do Kubernetes, os eventos de autenticação, elevação de privilégio e execução de comandos no sistema operacional dos nós e os fluxos das cargas de trabalho em contêineres. A inspeção dessa camada verifica a geração efetiva dos campos forenses obrigatórios em operação real.

A terceira fonte corresponde aos artefatos e eventos consolidados na plataforma de SIEM, contemplando decodificadores, regras de correlação mapeadas às exigências normativas, alertas indexados no repositório central e relatórios de avaliação automatizada de configuração do SIEM Wazuh \cite{wazuh2025}. Essa fonte permite auditar centralização, integridade, retenção e capacidade de detecção automatizada.

\todo[inline, color=pink]{A 1a. 2a. e 3a. fonte não são independentes e isso precisa ser explicado. É uma ameaça a validade que precisará ser tratada. Como elas não são independentes, a triangulação não é perfeita. Ainda assim, não indica quais achados serão corroborados por quais fontes, nem como divergências serão tratadas. Contudo, não existe triangulação metodológica. Todas as fontes são análise de artefatos (terceiro grau). Só o relatório automatizado do Wazuh se aproxima de coleta indireta (segundo grau). Verificar com cuidado. Eu não tenho muita clareza sobre isso.}

A quarta fonte abrange os documentos institucionais e os marcos normativos federais de referência, compostos por diretrizes operacionais internas de gestão de registros e pelos textos oficiais do Guia do Framework do PPSI 2.0 \cite{ppsi2026}, da Instrução Normativa SGD/MGI nº 4/2026 \cite{in4_2026} e da Norma Complementar nº 21/IN01/DSIC/GSI-PR \cite{nc21_2014}, que servem de gabarito objetivo para a avaliação.

\todo[inline, color=pink]{Os textos do PPSI 2.0, da IN SGD/MGI e da NC 21 não são dados sobre o caso, mas sim critérios normativos. A fonte de dados do caso são os artefatos técnicos e os registros de auditoria.}

\todo[inline, color=pink]{No refencial teórico, para as fontes de dados, você descreve primeiro, segundo e terceiro grau, mas aqui, você não explica. }

\section{Unidades de Análise e Procedimentos}
\label{sec:ec_unidades_procedimentos}

\todo[inline, color=pink]{Veja criteriosamente os apontamentos que fiz na introdução sobre as unidades de análise e aplique as mesmas correções aqui, mantendo a coerência com o restante do texto.}

No delineamento de estudo de caso único incorporado de \citeonline{yin2015} e \citeonline{runeson2009}, distingue-se formalmente o caso holístico, as subunidades de análise e os critérios de avaliação. O caso holístico corresponde à infraestrutura Kubernetes mantida pelo laboratório, enquanto as subunidades de análise compreendem os artefatos declarativos de configuração, os fluxos de registros em tempo de execução do plano de controle, dos nós e das cargas em contêineres, e os eventos, regras e relatórios consolidados na plataforma de SIEM. As doze medidas do Controle 8 do PPSI 2.0 \cite{ppsi2026} e os requisitos forenses da NC nº 21/GSI-PR \cite{nc21_2014} atuam como os critérios normativos contra os quais cada subunidade é avaliada.

A condução planejada para a segunda etapa da pesquisa divide-se em duas fases sequenciais. Na Fase 1, referente ao diagnóstico inicial de linha de base, executa-se coleta estritamente não intrusiva por análise arquival \cite{lethbridge2005,runeson2012} das configurações vigentes e dos registros disponíveis antes de qualquer modificação estrutural, preservando cada artefato com carimbo temporal e verificação de integridade. Na Fase 2, correspondente à intervenção arquitetural e à reavaliação, implanta-se uma arquitetura de monitoramento e auditoria de segurança em tempo de execução baseada na integração entre configurações nativas de auditoria e registro do Kubernetes, abrangendo plano de controle, nós e serviços internos, e a coleta, normalização, correlação e retenção centralizada pelo SIEM Wazuh, com vistas a viabilizar a verificação contínua da conformidade e elevar o Nível de Adoção do Controle 8 e o atendimento à NC nº 21/GSI-PR. Após a estabilização da solução, o roteiro de coleta da Fase 1 é integralmente reaplicado.

A Tabela~\ref{tab:ec_matriz_controle8} estabelece a matriz de verificação para todas as doze medidas do Controle 8, sem presumir previamente um Nível de Institucionalização fixo, especificando o requisito técnico e forense verificado, o mecanismo primário de evidência e a camada responsável pelo atendimento, distinguindo configuração nativa do Kubernetes e dos nós, plataforma SIEM Wazuh e processo organizacional.

\begin{table}[htbp]
    \centering
    \caption{Matriz de verificação das doze medidas do Controle 8 do PPSI 2.0 e da NC nº 21/GSI-PR}
    \label{tab:ec_matriz_controle8}
    \footnotesize
    \begin{tabular}{>{\centering\arraybackslash}p{0.7cm}>{\raggedright\arraybackslash}p{5.2cm}>{\raggedright\arraybackslash}p{5.2cm}>{\raggedright\arraybackslash}p{2.4cm}}
        \hline
        \textbf{ID} & \textbf{Requisito técnico e forense} & \textbf{Mecanismo primário de verificação e evidência} & \textbf{Camada responsável} \\ \hline
        8.1 & Processo documentado de gestão, coleta, revisão e retenção de registros com revisão periódica. & Inspeção de política interna de gestão de registros e artefatos de governança apoiados por relatórios do SIEM. & Processo e SIEM Wazuh \\ \hline
        8.2 & Coleta ativa de registros de autenticação, acesso a recursos e alterações de segurança nos ativos (NC 21, item 6.2). & Conferência de política de auditoria ativa no servidor de API, auditoria de nós e ingestão efetiva no índice central. & Kubernetes, nós e SIEM Wazuh \\ \hline
        8.3 & Armazenamento seguro com controle de acesso, capacidade adequada e integridade verificável (NC 21, item 7.5). & Verificação de permissões de acesso, cotas de disco, monitoramento de integridade e resumos criptográficos. & Nós e SIEM Wazuh \\ \hline
        8.4 & Sincronização temporal com ao menos duas fontes confiáveis alinhadas à Hora Legal Brasileira (NC 21, item 6.1). & Checagem da configuração e do estado do serviço NTP nos nós e consistência de carimbos temporais nos eventos. & Configuração nativa dos nós \\ \hline
        8.5 & Registros detalhados em ativos críticos contendo usuário, evento, data, hora, fuso e origem (NC 21, itens 6.3 e 6.5). & Auditoria dos níveis de verbosidade da política de auditoria do Kubernetes e validação dos quatro metadados nos eventos. & Kubernetes e SIEM Wazuh \\ \hline
        8.6 & Coleta de registros de consultas de resolução de nomes DNS no ambiente. & Verificação de diretivas de registro no serviço interno de DNS do \textit{cluster} e encaminhamento de eventos ao SIEM. & Kubernetes e SIEM Wazuh \\ \hline
        8.7 & Coleta de registros de requisições de URL em serviços expostos. & Verificação de registros de acesso HTTP nos componentes de roteamento de entrada e indexação no SIEM. & Kubernetes e SIEM Wazuh \\ \hline
        8.8 & Coleta de registros de execução em linha de comando nos hospedeiros. & Inspeção de regras de auditoria de chamadas e comandos no sistema operacional dos nós e alertas correlacionados. & Nós e SIEM Wazuh \\ \hline
        8.9 & Centralização e replicação remota da coleta e retenção de registros de auditoria (NC 21, item 6.8). & Comprovação de encaminhamento contínuo dos agentes para repositório central isolado fora dos nós de origem. & SIEM Wazuh \\ \hline
        8.10 & Retenção de registros pelo prazo mínimo normativo de seis meses, sem prejuízo de normas específicas (NC 21, item 6.7). & Inspeção de políticas de ciclo de vida de índices, rotação controlada e arquivamento auditável no repositório central. & SIEM Wazuh \\ \hline
        8.11 & Condução de revisões periódicas de anomalias em ciclo semanal ou inferior. & Definição de rotina documentada de auditoria apoiada por painéis, alertas priorizados e relatórios periódicos do SIEM. & Processo e SIEM Wazuh \\ \hline
        8.12 & Coleta de registros de provedores de serviços, plano de controle e aplicações críticas. & Validação da captura de trilhas do plano de controle Kubernetes e dos fluxos de saída das cargas em contêineres. & Kubernetes e SIEM Wazuh \\ \hline
    \end{tabular}
    \fonte{Elaborado pelo autor com base no PPSI 2.0 \cite{ppsi2026}, na IN SGD/MGI nº 4/2026 \cite{in4_2026} e na NC nº 21/GSI-PR \cite{nc21_2014}.}
\end{table}

\todo[inline, color=pink]{Essa tabela precisa estar na subsecao 4.5 estabelecendo como as medidas do Controle 8 são aplicadas nas unidades de análise.}


\section{Análise de Dados}
\label{sec:ec_analise_dados}

A análise de dados adota abordagem mista, segundo \citeonline{runeson2009} e \citeonline{wohlin2024}, combinando exame qualitativo técnico-normativo dos artefatos coletados e mensuração quantitativa determinística apoiada exclusivamente nos indicadores da IN SGD/MGI nº 4/2026 \cite{in4_2026}.

Na etapa qualitativa, aplica-se análise temática e de conteúdo sobre manifestos declarativos, registros brutos e relatórios do SIEM das Fases 1 e 2. Os artefatos são codificados segundo os requisitos da NC nº 21/GSI-PR \cite{nc21_2014}: sincronização com a Hora Legal Brasileira, registro de eventos críticos de autenticação e privilégio, presença dos quatro metadados mínimos obrigatórios, guarda pelo prazo normativo mínimo de seis meses, replicação remota e geração de resumos criptográficos de integridade. A triangulação evita que configurações declaradas, mas inoperantes, sejam computadas como conformes.

Na etapa quantitativa, as evidências trianguladas são convertidas de modo determinístico nas métricas oficiais da IN SGD/MGI nº 4/2026 \cite{in4_2026}. Para cada medida $i$ do Controle 8, atribui-se o indicador ordinal $\mathrm{NIM}_i \in \{0, 1, 2, 3, 4, 5\}$, cuja pontuação $p(\mathrm{NIM}_i)$ é dada pela Equação~\ref{eq:ec_nim}:
\begin{equation}
\label{eq:ec_nim}
p(\mathrm{NIM}_i) = \left\{
\begin{array}{ll}
0{,}00, & \mbox{se } \mathrm{NIM}_i = 1 \mbox{ (não adota)}, \\
0{,}25, & \mbox{se } \mathrm{NIM}_i = 2 \mbox{ (inicial, cobertura } < 40\%), \\
0{,}50, & \mbox{se } \mathrm{NIM}_i = 3 \mbox{ (intermediário, 40\% a 80\%)}, \\
0{,}75, & \mbox{se } \mathrm{NIM}_i = 4 \mbox{ (completo, cobertura } > 80\%), \\
1{,}00, & \mbox{se } \mathrm{NIM}_i = 5 \mbox{ (aprimorado, } > 80\% \mbox{ com revisão)}.
\end{array}
\right.
\end{equation}

Cada nível da Equação~\ref{eq:ec_nim} decorre da fração verificável de cobertura técnica e forense dos critérios da Tabela~\ref{tab:ec_matriz_controle8}, reservando-se o nível 5 aos casos com cobertura superior a 80\% associada a monitoramento contínuo e revisão quantitativa sistemática. O nível 0 indica não aplicabilidade justificada nos termos do parágrafo segundo do Art. 6º da IN SGD/MGI nº 4/2026 \cite{in4_2026}, admissível apenas quando o componente não existe no ambiente avaliado.

\todo[inline, color=pink]{Como ocorrerá o julgamento do nível de implementação da medida? Quem julgará? Qual o critério? Como será documentado? Essas questões são fundamentais para a validade do estudo de caso. É isso que ainda falta quando converso contigo sobre o como medir e interpretar.}

Consolidando as medidas aplicáveis, calcula-se o Nível de Adoção do Controle conforme a Equação~\ref{eq:ec_nac}:
\begin{equation}
\label{eq:ec_nac}
\mathrm{NAC} = \frac{1}{N_{\mathrm{aplic}}} \sum_{i \in M_{\mathrm{aplic}}} p(\mathrm{NIM}_i),
\end{equation}
em que $M_{\mathrm{aplic}}$ é o conjunto das medidas avaliadas e $N_{\mathrm{aplic}}$ sua cardinalidade, adotando-se integralmente as doze medidas como padrão, isto é, $N_{\mathrm{aplic}} = 12$, ressalvada dedução formalmente justificada conforme o parágrafo segundo do Art. 6º da norma \cite{in4_2026}. O índice enquadra o controle nas cinco faixas oficiais: nível 1 para $\mathrm{NAC} = 0{,}00$; nível 2 para $0{,}00 < \mathrm{NAC} < 0{,}25$; nível 3 para $0{,}25 \le \mathrm{NAC} < 0{,}50$; nível 4 para $0{,}50 \le \mathrm{NAC} < 0{,}75$; e nível 5 para $0{,}75 \le \mathrm{NAC} \le 1{,}00$.

O efeito da intervenção arquitetural entre a linha de base da Fase 1 e a reavaliação da Fase 2 é aferido pela variação absoluta da Equação~\ref{eq:ec_delta_nac}:
\begin{equation}
\label{eq:ec_delta_nac}
\Delta\mathrm{NAC} = \mathrm{NAC}_{\mathrm{final}} - \mathrm{NAC}_{\mathrm{inicial}},
\end{equation}
complementada pela análise desagregada $\Delta p(\mathrm{NIM}_i) = p(\mathrm{NIM}_{i,\mathrm{final}}) - p(\mathrm{NIM}_{i,\mathrm{inicial}})$ por medida, diferenciando ganhos diretos da arquitetura técnica daqueles dependentes de processos organizacionais.

\section{Instrumentação}
\label{sec:ec_instrumentacao}

A instrumentação define os recursos tecnológicos previstos para coleta, auditoria e centralização de evidências nas duas fases do estudo, sem vinculação a distribuições específicas ou ferramentas internas do laboratório.

Na camada de orquestração, utilizam-se subsistemas nativos do Kubernetes \cite{kubernetes2025}. O servidor de API \texttt{kube-apiserver} opera em conjunto com manifestos declarativos \texttt{AuditPolicy} para registrar requisições ao plano de controle segundo níveis graduais de detalhamento. O utilitário \texttt{kubectl} atua em modo estritamente de leitura para extrair manifestos e parâmetros de configuração durante a coleta não intrusiva, enquanto os componentes internos de resolução DNS e roteamento de entrada HTTP do \textit{cluster} subsidiam a avaliação das medidas 8.6 e 8.7.

Na camada dos nós hospedeiros Linux, empregam-se o subsistema nativo de auditoria do núcleo, voltado à captura de chamadas, comandos em linha de comando e alterações em arquivos críticos, e o serviço de sincronização temporal via protocolo \textit{Network Time Protocol} (NTP), calibrado a fontes redundantes vinculadas à Hora Legal Brasileira.

Na camada de monitoramento centralizado, adota-se o SIEM Wazuh \cite{wazuh2025}, composto por cinco elementos integrados. O \textit{Wazuh Agent} nos nós realiza coleta de registros e verificação de integridade de arquivos; o \textit{Wazuh Server}, denominado \textit{Manager}, executa decodificação e correlação por regras normativas; o \textit{Wazuh Indexer} garante armazenamento distribuído e retenção íntegra fora dos nós de origem; o \textit{Wazuh Dashboard} provê painéis e relatórios para auditoria periódica; e o módulo \textit{Security Configuration Assessment} executa verificações automatizadas de endurecimento frente a referenciais consolidados de segurança.


\todo[inline, color=pink]{Esta subseção precisa detalhar como cada ferramenta será configurada, quais parâmetros serão utilizados, quais logs serão coletados, quais regras de correlação serão aplicadas, como será feita a auditoria e a verificação de integridade, entre outros detalhes técnicos.}

\section{Ameaças à Validade}
\label{sec:ec_ameacas_validade}

Conforme a taxonomia de \citeonline{runeson2009} e \citeonline{yin2015}, estruturam-se as ameaças à validade do planejamento e as respectivas estratégias de mitigação para a condução na segunda etapa da pesquisa.

A \textbf{validade de constructo} enfrenta o risco de interpretação divergente entre os controles do PPSI 2.0, os requisitos forenses da NC nº 21/GSI-PR e as evidências extraídas do \textit{cluster}. Essa ameaça é mitigada pela matriz prévia de evidências por medida definida na Seção~\ref{sec:ec_unidades_procedimentos} e pelo emprego exclusivo dos indicadores oficiais NIM e NAC da IN SGD/MGI nº 4/2026 \cite{in4_2026}.

A \textbf{validade interna} pode ser afetada por alterações paralelas de configuração, atualizações de cargas de trabalho ou intervenções administrativas externas entre a linha de base e a reavaliação, interferindo na atribuição causal dos efeitos. Para mitigá-la, adotam-se o registro prévio do inventário do estado inicial, o controle rigoroso da janela temporal de coleta e o isolamento explícito dos manifestos e regras de auditoria introduzidos na Fase 2.

A \textbf{validade externa} decorre da impossibilidade de generalização estatística a partir de um estudo de caso único. Segundo \citeonline{yin2015}, essa limitação é tratada por meio da generalização analítica, descrevendo a infraestrutura e o protocolo em termos arquiteturais independentes de distribuição Kubernetes específica, assegurando transferibilidade a outros órgãos do SISP.

Por fim, a \textbf{confiabilidade} é ameaçada por eventual subjetividade na atribuição das notas NIM ou perda de rastreabilidade dos artefatos coletados. Como mitigação, formalizam-se os critérios determinísticos de pontuação na Seção~\ref{sec:ec_analise_dados}, o arquivamento íntegro e versionado das evidências brutas de cada fase com resumos criptográficos e a documentação auditável dos comandos e regras empregados.
